\apendice{Sostenibilización Curricular}
\label{apendice:F}

\section{Introducción}
El presente proyecto se ha desarrollado bajo la premisa de que la ingeniería informática debe contribuir de manera activa a la optimización de recursos y a la reducción del impacto ambiental. En este sentido, la plataforma de telerrehabilitación propuesta ofrece beneficios directos que van más allá del ámbito clínico, incidiendo en la eficiencia del sistema y en la sostenibilidad del entorno.

El aspecto más relevante de esta propuesta reside en la minimización sistemática de los desplazamientos en vehículos motorizados. Al permitir que la supervisión de los ejercicios se realice de forma remota y asíncrona, se elimina la necesidad de traslados recurrentes de pacientes y cuidadores hacia los centros de salud especializados. Esto se traduce en una reducción de la huella de carbono asociada al transporte y en un menor consumo de combustibles fósiles derivados de la logística sanitaria tradicional.

Desde el punto de vista social, el software mejora la accesibilidad del tratamiento al eliminar las barreras geográficas. Se garantiza que la calidad de la supervisión no dependa de la ubicación del usuario ni de su capacidad de desplazamiento físico, lo que democratiza el acceso a la rehabilitación. En definitiva, esta herramienta representa un avance hacia un modelo sanitario más eficiente y respetuoso con el medio ambiente, donde el uso de la tecnología permite optimizar la atención al paciente reduciendo los costes externos asociados.

\section{Alineación con los objetivos de desarrollo sostenible}

En concordancia con la Agenda 2030 de las Naciones Unidas\cite{un_agenda_2030}, el diseño e implementación de esta plataforma no solo responde a una necesidad funcional, sino que se alinea estratégicamente con varios objetivos de desarrollo sostenible. A continuación, se detalla la contribución del proyecto a los objetivos más relevantes:

\subsection{Eje social y sanitario}

\begin{description}
    \item[ODS 3: Salud y bienestar\cite{pactomundialODS3}.] \hfill \\
    La plataforma contribuye directamente a garantizar una vida sana y promover el bienestar. Al facilitar la telerrehabilitación asíncrona, se asegura la continuidad de los tratamientos y se incrementa la adherencia terapéutica. El sistema permite descongestionar las listas de espera presenciales, permitiendo que los recursos físicos se destinen a los casos más críticos, mientras que los pacientes en fases de mantenimiento reciben una supervisión constante y de calidad desde sus hogares.

    \item[ODS 10: Reducción de las desigualdades\cite{pactomundialODS10}.] \hfill \\
    Este proyecto incide en la reducción de la brecha de acceso a servicios especializados. Tradicionalmente, la calidad de la rehabilitación ha estado ligada a la proximidad a grandes núcleos urbanos con hospitales especializados. Esta herramienta elimina la barrera geográfica, ofreciendo equidad en el tratamiento a pacientes de zonas rurales (la España vaciada) o a personas con movilidad reducida severa que dependen de terceros para su traslado.
\end{description}

\subsection{Eje medioambiental y tecnológico}

\begin{description}
    \item[ODS 13: Acción por el clima\cite{pactomundialODS13}.] \hfill \\
    La contribución más cuantificable del proyecto es la reducción de emisiones de gases de efecto invernadero. Al sustituir el transporte físico de pacientes por el transporte digital de datos, se mitiga la huella de carbono asociada a la movilidad urbana e interurbana. Esto apoya las políticas de descarbonización y promueve una gestión más eficiente de la energía en el sector servicios.

    \item[ODS 9: Industria, innovación e infraestructura\cite{pactomundialODS9}.] \hfill \\
    El software desarrollado es un ejemplo de modernización de la infraestructura sanitaria mediante la digitalización. Se optimiza el uso de recursos técnicos y humanos, demostrando cómo la ingeniería informática puede transformar procesos tradicionales en sistemas más resilientes y sostenibles.
\end{description}



\section{Conclusión}
En conclusión, la plataforma presentada trasciende la mera solución técnica para constituirse como una herramienta de transformación social y ecológica. La integración de criterios de sostenibilidad en el ciclo de vida del software ha permitido desarrollar un producto que no solo cura, sino que cuida el entorno y democratiza el acceso a la salud, cumpliendo con la responsabilidad ética inherente a la ingeniería contemporánea.